% GENERATED FILE. Edit canonical lessons, then run npm run generate:books.
% canonical-source-hash: fnv1a64:f0ea9cfa627bbb27
% canonical-lessons: GU-C173-athvadiyu, GU-C173-pakhvadiyu, GU-C173-avtu, GU-C173-chellu, GU-C173-akhu

\chapter{A week and a fortnight}
\label{ch:gu-a-week-and-a-fortnight}
\hlchaptermodality{173}

\begin{chapteropening}
\textbf{By the end of this chapter:} \emph{I can say how long I will stay and what the last day will be.}

Complete the last lesson of chapter 173: I can say how long I will stay and what the last day will be.
\end{chapteropening}

\section[\texorpdfstring{\gu{અઠવાડિયું} (aṭhvāḍiyũ)}{અઠવાડિયું (aṭhvāḍiyũ)}]{\gu{અઠવાડિયું} (aṭhvāḍiyũ) — week}

\label{lesson:GU-C173-athvadiyu}



\begin{quote}
Before the new one: say the Gujarati for orange, then the Gujarati for juice.
\end{quote}

\subsection*{You'll want to know: \gu{અઠવાડિયું}}
\textbf{\gu{અઠવાડિયું}} — \emph{aṭhvāḍiyũ} — \textquotedblleft{}week\textquotedblright{}.

\begin{quote}
\textbf{\gu{હું} \gu{એક} \gu{અઠવાડિયું} \gu{ગામમાં} \gu{રહીશ}.} — \emph{hũ ek aṭhvāḍiyũ gāmmā̃ rahīś} — \textquotedblleft{}I will stay in the village for a week.\textquotedblright{}
\end{quote}

\begin{cousinweb}
\textbf{\gu{અઠવાડિયું}} (\emph{aṭhvāḍiyũ}) is a tadbhava. It comes from \emph{*aṣṭavāra}, \textquotedblleft{}period of eight days\textquotedblright{}, a reconstructed form built from Sanskrit \emph{aṣṭā}, \textquotedblleft{}eight\textquotedblright{}, and \emph{vāra}, \textquotedblleft{}day\textquotedblright{}; CDIAL lists the Gujarati \emph{aṭhvāḍiyũ} and takes it to be crossed with compounds of \emph{vāṭa}, \textquotedblleft{}enclosure\textquotedblright{} (CDIAL 940).
\end{cousinweb}

\begin{grammarlens}[title={What you've built}]
One more word for everyday life at A2.
\end{grammarlens}

\subsection*{Your turn}
\begin{itemize}
\raggedright
  \item \emph{Say it:} \emph{aṭhvāḍiyũ}
  \item \emph{Say it:} \emph{aṭhvāḍiyũ}, once more
  \item \emph{Say it:} \emph{hũ ek aṭhvāḍiyũ gāmmā̃ rahīś} — I will stay in the village for a week.
  \item \emph{Recall:} say the Gujarati for orange, then the Gujarati for juice, then say \emph{aṭhvāḍiyũ} again
\end{itemize}

\begin{tcolorbox}[breakable,colback=teal!4,colframe=teal!35!black,title={Before you move on}]
What does \gu{અઠવાડિયું} mean? (\textquotedblleft{}Week\textquotedblright{}.) Say it once more.
\end{tcolorbox}
\section[\texorpdfstring{\gu{પખવાડિયું} (pakhvāḍiyũ)}{પખવાડિયું (pakhvāḍiyũ)}]{\gu{પખવાડિયું} (pakhvāḍiyũ) — fortnight}

\label{lesson:GU-C173-pakhvadiyu}



\begin{quote}
Before the new one: say the Gujarati for juice, then the Gujarati for week.
\end{quote}

\subsection*{You'll want to know: \gu{પખવાડિયું}}
\textbf{\gu{પખવાડિયું}} — \emph{pakhvāḍiyũ} — \textquotedblleft{}fortnight\textquotedblright{}.

\begin{quote}
\textbf{\gu{તે} \gu{એક} \gu{પખવાડિયું} \gu{રજા} \gu{લેશે}.} — \emph{te ek pakhvāḍiyũ rajā leśe} — \textquotedblleft{}She will take a fortnight off.\textquotedblright{}
\end{quote}

\begin{cousinweb}
\textbf{\gu{પખવાડિયું}} (\emph{pakhvāḍiyũ}) is built like \emph{aṭhvāḍiyũ}, \textquotedblleft{}week\textquotedblright{}, on \emph{pakh}, \textquotedblleft{}side, half\textquotedblright{}. CDIAL lists Gujarati \emph{pākh}, \textquotedblleft{}side\textquotedblright{}, under Sanskrit \emph{pakṣa}, \textquotedblleft{}wing, side\textquotedblright{}, but not this compound, so its path is not settled (CDIAL 7627).
\end{cousinweb}

\begin{grammarlens}[title={What you've built}]
One more word for everyday life at A2.
\end{grammarlens}

\subsection*{Your turn}
\begin{itemize}
\raggedright
  \item \emph{Say it:} \emph{pakhvāḍiyũ}
  \item \emph{Say it:} \emph{pakhvāḍiyũ}, once more
  \item \emph{Say it:} \emph{te ek pakhvāḍiyũ rajā leśe} — She will take a fortnight off.
  \item \emph{Recall:} say the Gujarati for juice, then the Gujarati for week, then say \emph{pakhvāḍiyũ} again
\end{itemize}

\begin{tcolorbox}[breakable,colback=teal!4,colframe=teal!35!black,title={Before you move on}]
What does \gu{પખવાડિયું} mean? (\textquotedblleft{}Fortnight\textquotedblright{}.) Say it once more.
\end{tcolorbox}
\section[\texorpdfstring{\gu{આવતું} (āvtũ)}{આવતું (āvtũ)}]{\gu{આવતું} (āvtũ) — next, coming}

\label{lesson:GU-C173-avtu}



\begin{quote}
Before the new one: say the Gujarati for week, then the Gujarati for fortnight.
\end{quote}

\subsection*{You'll want to know: \gu{આવતું}}
\textbf{\gu{આવતું}} — \emph{āvtũ} — \textquotedblleft{}next, coming\textquotedblright{}.

\begin{quote}
\textbf{\gu{આવતા} \gu{અઠવાડિયે} \gu{હું} \gu{આવીશ}.} — \emph{āvtā aṭhvāḍiye hũ āvīś} — \textquotedblleft{}I will come next week.\textquotedblright{}
\end{quote}

Built on \textbf{\gu{આવવું}}, \emph{āvvũ}, to come.

\begin{cousinweb}
\textbf{\gu{આવતું}} (\emph{āvtũ}) is built on the inherited \emph{āvvũ}, \textquotedblleft{}to come\textquotedblright{}, with the participle ending \emph{-tũ}, so its first sense is \textquotedblleft{}coming\textquotedblright{}.
\end{cousinweb}

\begin{grammarlens}[title={What you've built}]
One more word for everyday life at A2.
\end{grammarlens}

\subsection*{Your turn}
\begin{itemize}
\raggedright
  \item \emph{Say it:} \emph{āvtũ}
  \item \emph{Say it:} \emph{āvtũ}, once more
  \item \emph{Say it:} \emph{āvtā aṭhvāḍiye hũ āvīś} — I will come next week.
  \item \emph{Recall:} say the Gujarati for week, then the Gujarati for fortnight, then say \emph{āvtũ} again
\end{itemize}

\begin{tcolorbox}[breakable,colback=teal!4,colframe=teal!35!black,title={Before you move on}]
What does \gu{આવતું} mean? (\textquotedblleft{}Next, coming\textquotedblright{}.) Say it once more.
\end{tcolorbox}
\section[\texorpdfstring{\gu{છેલ્લું} (chellũ)}{છેલ્લું (chellũ)}]{\gu{છેલ્લું} (chellũ) — last}

\label{lesson:GU-C173-chellu}



\begin{quote}
Before the new one: say the Gujarati for fortnight, then the Gujarati for next.
\end{quote}

\subsection*{You'll want to know: \gu{છેલ્લું}}
\textbf{\gu{છેલ્લું}} — \emph{chellũ} — \textquotedblleft{}last\textquotedblright{}.

\begin{quote}
\textbf{\gu{આ} \gu{છેલ્લો} \gu{દિવસ} \gu{હશે}.} — \emph{ā chello divas haśe} — \textquotedblleft{}This will be the last day.\textquotedblright{}
\end{quote}

\begin{cousinweb}
The origin of \textbf{\gu{છેલ્લું}} (\emph{chellũ}) is not settled. CDIAL lists no Gujarati form of it, and EtymDB has no entry for it.
\end{cousinweb}

\begin{grammarlens}[title={What you've built}]
One more word for everyday life at A2.
\end{grammarlens}

\subsection*{Your turn}
\begin{itemize}
\raggedright
  \item \emph{Say it:} \emph{chellũ}
  \item \emph{Say it:} \emph{chellũ}, once more
  \item \emph{Say it:} \emph{ā chello divas haśe} — This will be the last day.
  \item \emph{Recall:} say the Gujarati for fortnight, then the Gujarati for next, then say \emph{chellũ} again
\end{itemize}

\begin{tcolorbox}[breakable,colback=teal!4,colframe=teal!35!black,title={Before you move on}]
What does \gu{છેલ્લું} mean? (\textquotedblleft{}Last\textquotedblright{}.) Say it once more.
\end{tcolorbox}
\section[\texorpdfstring{\gu{આખું} (ākhũ)}{આખું (ākhũ)}]{\gu{આખું} (ākhũ) — whole}

\label{lesson:GU-C173-akhu}



\begin{quote}
Before the new one: say the Gujarati for next, then the Gujarati for last.
\end{quote}

\subsection*{You'll want to know: \gu{આખું}}
\textbf{\gu{આખું}} — \emph{ākhũ} — \textquotedblleft{}whole\textquotedblright{}.

\begin{quote}
\textbf{\gu{હું} \gu{આખો} \gu{દિવસ} \gu{કામ} \gu{કરીશ}.} — \emph{hũ ākho divas kām karīś} — \textquotedblleft{}I will work the whole day.\textquotedblright{}
\end{quote}

\begin{cousinweb}
\textbf{\gu{આખું}} (\emph{ākhũ}) is a tadbhava. It comes from Sanskrit \emph{akṣata}, \textquotedblleft{}unhurt, not crushed\textquotedblright{}, and CDIAL glosses the Gujarati word \textquotedblleft{}whole\textquotedblright{} (CDIAL 28).
\end{cousinweb}

\begin{grammarlens}[title={What you've built}]
One more word for everyday life at A2.
\end{grammarlens}

\subsection*{Your turn}
\begin{itemize}
\raggedright
  \item \emph{Say it:} \emph{ākhũ}
  \item \emph{Say it:} \emph{ākhũ}, once more
  \item \emph{Say it:} \emph{hũ ākho divas kām karīś} — I will work the whole day.
  \item \emph{Recall:} say the Gujarati for next, then the Gujarati for last, then say \emph{ākhũ} again
\end{itemize}

\begin{tcolorbox}[breakable,colback=teal!4,colframe=teal!35!black,title={Before you move on}]
What does \gu{આખું} mean? (\textquotedblleft{}Whole\textquotedblright{}.) Say it once more.
\end{tcolorbox}
